% !PRINTED 226
\weberpage{226}
Nun ist nach \S.\,68, (6)
\[
S f_0 = n a_0, \quad S f_1 = (n-1) a_1, \ldots S f_{n-1} = a_{n-1},
\]
so dass man erhält
\[
\begin{aligned}
S(y z) = {} & n a_0 \sum_{0,\,n-1}^{s} E_{0,s} \tau_{n-s-1} \\
& + (n-1) a_1 \sum_{0,\,n-1}^{s} E_{1,s} \tau_{n-s-1} \\
& \cdots\cdots\cdots\cdots \\
& + a_{n-1} \sum_{0,\,n-1}^{s} E_{n-1,s} \tau_{n-s-1}.
\end{aligned}
\]
Hierin ist nun der Coëfficient von $\tau_{n-1}$
\[
n a_0 E_{0,0} + (n-1) a_1 E_{1,0} + \cdots + a_{n-1} E_{n-1,0},
\]
also nach (3)
\[
= a_0 S(y).
\]
Da nun $S(z)$ mit dem Gliede $n a_0 \tau_{n-1}$ anfängt, so wird in
der Differenz
\weq{5}{S(y z) - \frac{1}{n} S(y) S(z)}
kein Glied vorkommen, das mit $\tau_{n-1}$ multiplicirt ist. Da dieser
Ausdruck aber ausserdem in Bezug auf $t$ und $\tau$ symmetrisch ist,
so enthält er auch nicht $t_{n-1}$, und wir können bei seiner Bildung
% [sic] Der Druck lässt nach dem ersten Glied „=0“ aus; der unmittelbar folgende Absatz bestätigt, dass sowohl t_{n-1} als auch tau_{n-1} gleich Null gesetzt werden.
einfach $t_{n-1}$ und $\tau_{n-1} = 0$ annehmen.

Wir setzen nun
\[
S(y) = n a_0 t_{n-1} + \varphi(t),
\]
worin
\weq{6}{\varphi(t) = (n-1) a_1 t_{n-2} + (n-2) a_2 t_{n-3} + \cdots + a_{n-1} t_0.}
Dann wird
\weq{7}{
\begin{aligned}
S(y z) - \frac{1}{n} S(y) S(z) = {} & n a_0 \sum_{1,\,n-1}^{s} E_{0,s} \tau_{n-s-1} \\
& + (n-1) a_1 \sum_{1,\,n-1}^{s} E_{1,s} \tau_{n-s-1} \\
& \cdots\cdots\cdots\cdots \\
& + a_{n-1} \sum_{1,\,n-1}^{s} E_{n-1,s} \tau_{n-s-1} - \frac{1}{n} \varphi(t) \varphi(\tau),
\end{aligned}
}
wobei jedoch zu bemerken ist, dass $t_{n-1} = 0$ anzunehmen, also
\[
E_{s,s} = a_1 t_{n-2} + \cdots + a_s t_{n-s-1}
\]
zu setzen ist, während die übrigen $E$ durch die Formeln (9) des
vorigen Paragraphen bestimmt sind.
